% SHARD-ID: CA-30-GAUSSIAN-BOSON-1D-STRESS-CANDIDATES
% SHARD-TITLE: One-Dimensional Gaussian Boson Stress-Energy Candidates
% SHARD-SUMMARY: Proposes a 1+1-dimensional lattice stress-energy dictionary for the scalar Gaussian chain using the CA-29 density split.
% SHARD-SUMMARY: Promotes only the finite open-chain energy-current component T_01 to a checked seed via the A4 continuity witness.
% SHARD-SUMMARY: Keeps T_10 and T_11 as future momentum-density and momentum-current targets, not as consequences of symmetry or tracelessness.
% SHARD-KEYWORDS: Gaussian boson, stress-energy candidate, energy current, momentum density, open chain, continuity equation

\section{One-Dimensional Gaussian Boson Stress-Energy Candidates}
\label{sec:gaussian-boson-1d-stress-candidates}

\subsection{Status, Sources, and Dependencies}

\textbf{Status: proposal with checked \(T_{01}\) seed.}  This shard fixes a
candidate dictionary for the finite open one-dimensional scalar Gaussian chain.
Only the energy-current entry has a checked local witness.  No convergence
theorem, no continuum stress tensor, and no Virasoro representation is claimed.

Dependencies: CA-12, CA-23--CA-29, CONVENTIONS.md (j), (k), (l), and the A4
finite open-chain current tests.

% Source: references/qft/SOURCES.md:59--90 -- registered Weinberg A1 anchor ranges and improvement-status boundary.
% Source: references/qft/SOURCES.md:94--97 -- targeted OCR is an anchor convenience; exact signs/indices need page-image checks.
% Source: references/qft/Weinberg1995/Weinberg1995_QFT1_A1_ocr.txt:260--280 -- translations give conserved currents grouped as the energy-momentum tensor and integrated translation generators.
% Source: references/qft/Weinberg1995/Weinberg1995_QFT1_A1_ocr.txt:357--361 -- canonical raised tensor need not be symmetric.
% Source: references/qft/Weinberg1995/Weinberg1995_QFT1_A1_ocr.txt:572--589 -- Belinfante tensor is conserved, has the same integrated translations, and is symmetric.
% Source: references/qft/Weinberg1995/Weinberg1995_QFT1_A1_ocr.txt:619--633 -- boost generator and displayed commutators with H and P.
% Source: references/qft/Weinberg1995/Weinberg1995_QFT1_A1_ocr.txt:675--684 -- Lorentz commutator proof caveat for boost generators with auxiliary fields.
% Source: references/cft/Schottenloher2008/Schottenloher2008.md:5046--5052 -- Axiom 5: 1+1 CFT energy-momentum tensor fields, symmetry, and conservation.
% Source: references/cft/Schottenloher2008/Schottenloher2008.md:5054--5063 -- Luescher-Mack consequence: tracelessness and Virasoro modes under Axioms 1--5.
% Convention: CONVENTIONS.md:219--239 -- finite periodic checks do not define first moments or coordinate generators without a branch policy.
% Convention: CONVENTIONS.md:312--347 -- Poincare vector-field brackets translated to self-adjoint commutator signs.
% Convention: CONVENTIONS.md:349--421 -- Gaussian density split, open-chain current orientation, endpoint convention, subtraction, and improvement policies.
% Source: report/sections/12_lattice_boost_current_1d.tex:24--89 -- first energy moment yields nearest-neighbour current candidate with periodic position warning.
% Source: report/sections/29_gaussian_boson_real_space_density.tex:26--42 -- physical density h_x versus integrated site energy e_x.
% Source: report/sections/29_gaussian_boson_real_space_density.tex:93--105 -- open-boundary and periodic first-moment policies.
% Source: report/sections/29_gaussian_boson_real_space_density.tex:107--151 -- subtraction and improvement policies.
% Source: src/GaussianBosonCurrents.jl:127--183 -- finite open-chain energy-current and continuity-residual helpers.
% Checked: test/runtests.jl:236--267 -- "Gaussian open-chain energy current continuity".

\subsection{Continuum Motivation and Boundary}

The continuum notation \(T_{\mu\nu}\) is used here as a target shape, not as an
imported lattice theorem.  Weinberg supplies the QFT motivation that spacetime
translations lead to conserved currents grouped as an energy-momentum tensor and
that translation generators are spatial integrals of time components.  The same
registered OCR also records two caveats that matter for the lattice dictionary:
the canonical raised tensor need not be symmetric, while the Belinfante tensor is
a modified conserved tensor with the same integrated translation generators and
symmetry.  The OCR around formulas is noisy, so this shard uses those anchors
only for motivation and caveats, not for exact component signs.
The self-adjoint commutator signs for any later Poincare comparison are the
project convention in CONVENTIONS.md (k), not raw OCR import.

Schottenloher's Axiom 5 is narrower but sharper for the CFT endpoint: in a
two-dimensional conformal field theory there are four fields \(T_{\mu\nu}\),
\(\mu,\nu\in\{0,1\}\), with symmetry and conservation.  The later
Luescher--Mack theorem gives tracelessness and two Virasoro mode families under
the full Axioms 1--5.  We do not use those CFT conclusions to define the lattice
\(T_{10}\) or \(T_{11}\).  They are acceptance targets for a future continuum
proof.

\subsection{Finite Open Chain Setup}

Work first with a finite open chain of \(N\) cells and spacing \(\epsilon\).  The
CA-29 density convention distinguishes the physical density \(h_j\) from the
integrated cell energy
\[
  e_j=\epsilon h_j,
  \qquad
  H_\epsilon=\sum_{j=1}^N e_j .
\]
Finite commutators use \(e_j\).  This avoids inserting an accidental factor of
\(\epsilon\) into a discrete continuity identity.

For the nearest-neighbour scalar Gaussian chain, the A4 helper implements the
CA-29 equal-endpoint split of the finite Hamiltonian: kinetic and onsite terms
belong to the site, and each present bond \(q_j q_{j+1}\) is split equally
between its two endpoints.  Open-chain means no exterior ghost cells, no missing
half-bonds, and no boundary counterterms are implicit.

\subsection{Candidate Dictionary}

\paragraph{\(T_{00}^{\mathrm{lat}}\).}
The lattice energy-density entry is the CA-29 density:
\[
  T_{00}^{\mathrm{lat}}(j):=h_j .
\]
For exact finite open-chain algebra, the cell-integrated version is
\[
  E_j:=e_j=\epsilon h_j .
\]
Thus \(T_{00}^{\mathrm{lat}}\) is the physical density label, while \(E_j\) is
the object that appears in the tested quadratic commutators.

\paragraph{\(T_{01}^{\mathrm{lat}}\).}
The energy-current or energy-flux seed lives on open-chain edges:
\[
  J_{j+1/2}:=i[e_j,e_{j+1}],
  \qquad 1\le j<N .
\]
Set endpoint currents to zero,
\[
  J_{1/2}=J_{N+1/2}=0 .
\]
The checked finite continuity identity is
\[
  i[H_\epsilon,e_j]
  =
  J_{j-1/2}-J_{j+1/2},
  \qquad 1\le j\le N .
\]
This is the only stress-energy component in this shard with a Julia witness.
It is the lattice candidate corresponding to \(T_{01}\) in the cell-integrated
continuity law
\[
  \frac{d}{dt}\int_{\text{cell }j}T_{00}\,dx
  =
  T_{01}(x_{j-1/2})-T_{01}(x_{j+1/2}) .
\]
No continuum normalization or convergence of \(J_{j+1/2}\) is claimed.

\paragraph{\(T_{10}^{\mathrm{lat}}\).}
Do not identify \(T_{10}^{\mathrm{lat}}\) with \(T_{01}^{\mathrm{lat}}\) at this
stage.  Define it instead as the future integrated momentum-density comparison
target:
\[
  M_j \quad\text{with}\quad P^{\mathrm{lat}}=\sum_j M_j .
\]
The required next witness is that the chosen \(M_j\), or at least its total
\(P^{\mathrm{lat}}\), matches the Gaussian one-particle flat-symbol momentum
target after the appropriate \(\mathrm{d}\Gamma(k)\) comparison is made.  That is
the S3 task, not an assumption in this shard.

\paragraph{\(T_{11}^{\mathrm{lat}}\).}
Define \(T_{11}^{\mathrm{lat}}\) operationally as the future current of the
chosen momentum density.  Once \(M_j\) is fixed, the candidate stress component
must come with edge terms \(S_{j+1/2}\) and a checked equation
\[
  i[H_\epsilon,M_j]
  =
  S_{j-1/2}-S_{j+1/2}
\]
with an explicit boundary convention.  Until that witness exists, there is no
lattice \(T_{11}\) formula.  In particular, \(T_{11}^{\mathrm{lat}}\) is not
defined as \(-T_{00}^{\mathrm{lat}}\), not derived from tracelessness, and not
inferred from CFT holomorphic factorization.

\subsection{Boundary, Subtraction, and Improvement}

This candidate dictionary is open-boundary first.  Endpoint currents are part of
the statement, not small errors:
\[
  i[H_\epsilon,e_1]=-J_{3/2},
  \qquad
  i[H_\epsilon,e_N]=J_{N-1/2}.
\]
Periodic total-energy density splits are allowed by CA-29, but periodic first
moments and periodic coordinate-dependent generator tests remain blocked until
CONVENTIONS.md (j) fixes a coordinate branch, sawtooth, Fourier interpolation,
or equivalent policy.

The default density is bare.  A state-dependent subtraction
\({:}h_j{:}_\omega=h_j-\omega(h_j)\mathbf 1\) may later be named; it changes the
Hamiltonian by a scalar and does not change commutators, but it is part of any
claim about absolute energy or density modes.  The default improvement is
zero.  Any later replacement \(h'_j=h_j+(\nabla_\epsilon b)_j\) must name the
edge field, discrete divergence, and boundary convention before first moments or
currents are recomputed.

\subsection{Compiler Payload and Acceptance}

A compiler output that cites this shard must carry:
\[
\begin{array}{ll}
\text{density convention} & h_j,\ e_j=\epsilon h_j,\ \text{CA-29 split},\\
\text{boundary policy} & \text{finite open chain, }J_{1/2}=J_{N+1/2}=0,\\
\text{subtraction policy} & \text{bare unless a state }\omega\text{ is named},\\
\text{improvement policy} & b=0\text{ unless a discrete divergence is named},\\
\text{checked }T_{01}\text{ witness} &
  i[H_\epsilon,e_j]=J_{j-1/2}-J_{j+1/2}.
\end{array}
\]

The acceptance boundary is equally explicit.  The checked witness is the finite
open-chain continuity residual in \texttt{src/GaussianBosonCurrents.jl} and
\texttt{test/runtests.jl}.  Still pending are the S3 equivalence witness
comparing \(T_{01}\) current or first-moment data with the flat
\(\mathrm{d}\Gamma(k)\) momentum target; the \(T_{11}\) witness consisting of a
chosen momentum density \(M_j\) and checked current \(S_{j+1/2}\); and the
continuum witness: scaling maps, domains, state sequence, and covariance
convergence.
Without those witnesses, the dictionary remains a proposal with one checked
component, not a stress-tensor theorem.
